% Parte compartida: paquetes, tipografía, colores y macros de diseño.
% Compilar con LuaLaTeX (ver build.sh).
\usepackage[a4paper,top=22mm,bottom=24mm,left=24mm,right=24mm]{geometry}
% Compila con pdfLaTeX (lo que usa Overleaf, más rápido) o con LuaLaTeX.
\usepackage{iftex}
\ifLuaTeX
  \usepackage[sfdefault]{FiraSans}   % misma tipografía que la marca (Google Fonts: Fira Sans)
  \setmonofont{FiraMono-Regular.otf}[BoldFont=FiraMono-Bold.otf,ItalicFont=FiraMono-Oblique.otf,Scale=MatchLowercase] % sin ligaduras TeX: comillas rectas
\else
  \usepackage[T1]{fontenc}
  \usepackage[utf8]{inputenc}
  \usepackage[sfdefault]{FiraSans}
  \usepackage{FiraMono}
  \DeclareUnicodeCharacter{2192}{\ensuremath{\rightarrow}}
\fi
\usepackage[spanish,es-noquoting,es-noshorthands]{babel}
\usepackage{tikz}
\usepackage{xcolor,graphicx,tabularx,array,enumitem,booktabs,colortbl}
\usepackage{fancyvrb,titlesec,fancyhdr,needspace}
\usepackage[most]{tcolorbox}
\usepackage[hidelinks]{hyperref}

% --- Colores de roles (iguales a diagram.svg) ---
\definecolor{ink}{HTML}{0F172A}
\definecolor{muted}{HTML}{475569}
\definecolor{line}{HTML}{CBD5E1}
\definecolor{soft}{HTML}{F1F5F9}
\definecolor{datos}{HTML}{EA580C}
\definecolor{server}{HTML}{16A34A}
\definecolor{algo}{HTML}{2563EB}
\definecolor{pitch}{HTML}{9333EA}
% --- Señales del guardián (separadas de la marca) ---
\definecolor{amber}{HTML}{CA8A04}
\definecolor{block}{HTML}{DC2626}
\definecolor{allow}{HTML}{16A34A}
\definecolor{codecol}{HTML}{0F5F6E}  % color del texto de código en este documento

% --- Texto con aire: una columna, interlineado amplio, párrafos separados ---
\color{ink}
\linespread{1.18}
\setlength{\parindent}{0pt}
\setlength{\parskip}{7pt}
\AtBeginDocument{\raggedright}
\setlist{itemsep=4pt,topsep=4pt,leftmargin=18pt}
\renewcommand{\arraystretch}{1.45}
\setlength{\tabcolsep}{6pt}
\arrayrulecolor{line}

% --- Pie de página ---
\pagestyle{fancy}
\fancyhf{}
\renewcommand{\headrulewidth}{0pt}
\fancyfoot[L]{\small\color{muted}\Project}
\fancyfoot[R]{\small\color{muted}\thepage}

% --- Partes y secciones ---
% \docpart{color}{título}{subtítulo}: empieza una parte en página nueva.
\newcommand{\partcolor}{ink}
\newcommand{\partlabel}{Parte}
\newcommand{\docpart}[3]{%
  \clearpage\refstepcounter{part}\setcounter{section}{0}%
  \renewcommand{\partcolor}{#1}%
  \addcontentsline{toc}{part}{\partlabel{} \thepart\quad #2}%
  {\color{#1}\rule{\linewidth}{4pt}}\par\vspace{6pt}
  {\color{#1}\bfseries\large\MakeUppercase{\partlabel} \thepart}\par\vspace{2pt}
  {\bfseries\fontsize{24}{28}\selectfont\hyphenpenalty=10000 #2\par}\vspace{6pt}
  {\color{muted}\large #3}\par\vspace{14pt}}
\renewcommand{\thepart}{\arabic{part}}
\renewcommand{\thesection}{\thepart.\arabic{section}}
\renewcommand{\theHsection}{\thepart.\arabic{section}} % anclas únicas para los enlaces del índice
\titleformat{\section}{\needspace{6\baselineskip}\Large\bfseries}{{\color{\partcolor}\thesection}}{0.6em}{}[{\color{line}\titlerule[0.8pt]}]
\titlespacing*{\section}{0pt}{20pt}{10pt}
\titleformat{\subsection}{\large\bfseries}{}{0pt}{}
\titlespacing*{\subsection}{0pt}{14pt}{4pt}
\setcounter{secnumdepth}{1}
\setcounter{tocdepth}{1}
\makeatletter\renewcommand{\l@section}{\@dottedtocline{1}{1.2em}{2.6em}}\makeatother

% --- Macros ---
\newcommand{\code}[1]{{\leavevmode\ttfamily\color{codecol}#1}}
\newcommand{\roletag}[2]{{\color{#1}\rule[-1pt]{4pt}{10pt}}\,\textbf{#2}}
\newcommand{\verdict}[2]{\tcbox[on line,colback=#1,colframe=#1,boxrule=0pt,arc=2pt,left=3pt,right=3pt,top=1pt,bottom=1pt]{\color{white}\bfseries\small\ttfamily #2}}
\newcommand{\Allow}{\verdict{allow}{allow}}
\newcommand{\BlockAttack}{\verdict{block}{block\_attack}}
\newcommand{\BlockTopic}{\verdict{block}{block\_off\_topic}}

\newtcolorbox{callout}[1][ink]{enhanced,boxrule=0pt,leftrule=3pt,colframe=#1,colback=soft,
  arc=0pt,left=10pt,right=10pt,top=8pt,bottom=8pt,before skip=10pt,after skip=12pt,
  before upper={\raggedright\setlength{\parskip}{6pt}}}
% Tarjeta con barra de color: \begin{card}{color}{título} ... \end{card}
\newtcolorbox{card}[2]{enhanced,boxrule=0.6pt,leftrule=5pt,colframe=#1,colback=white,
  arc=3pt,left=12pt,right=12pt,top=10pt,bottom=8pt,before skip=14pt,after skip=12pt,
  before upper={\raggedright\setlength{\parskip}{6pt}},
  fonttitle=\bfseries\large,coltitle=#1,colbacktitle=white,title={#2},
  attach boxed title to top left={yshift=-2mm,xshift=8pt},boxed title style={boxrule=0pt,colframe=white}}
% Bloques de código: \begin{codebox}\begin{Verbatim}...\end{Verbatim}\end{codebox}
\fvset{fontsize=\small}
\newtcolorbox{codebox}{colback=soft,colframe=line,boxrule=0.5pt,arc=2pt,left=8pt,right=8pt,top=4pt,bottom=4pt,before skip=8pt,after skip=10pt}
